% !TEX root = ../main.tex

\clearpage

\section{Conclusion and Perspectives}

\subsection{General Conclusion}
This internship project focused on the design and implementation of VIP -- Value Intelligent Platform, an enterprise-grade Agentic AI Hub for building, deploying, orchestrating, and monitoring intelligent applications at scale. The main objective was to address the limitations of fragmented AI solutions by proposing a modular platform where AI agents can be created from declarative configurations, deployed as independent services, composed into workflows, and monitored through observability mechanisms.

The work carried out on the VIP platform resulted in a configurable and extensible architecture based on several core principles. First, the platform follows a configuration-driven approach, where the behavior of an agent is defined through prompts, model parameters, tools, and input/output schemas rather than through direct source code modifications. This supports the ``Zero Code Change'' philosophy and makes it possible to create specialized agents while keeping the runtime logic generic and reusable.

Second, the platform introduced an Agent Factory capable of generating runnable agent services from user-defined configurations. Each generated agent is packaged as an independent FastAPI service and deployed inside a Docker container. This approach improves modularity, separation of concerns, and deployment flexibility. The generated agents expose a standardized runtime contract, which allows the orchestrator to invoke them consistently regardless of their internal role.

Third, the project implemented a workflow orchestration layer responsible for coordinating multiple agents in a structured execution pipeline. In the implemented financial document intelligence scenario, the workflow was composed of four main agents: an extractor agent, a mapper agent, a generator agent, and a sanitizer agent. The orchestrator ensured that each step received the appropriate input, validated payloads against schemas, executed agents in sequence, and returned a normalized workflow result.

A major contribution of the project was also the integration of governance and traceability mechanisms. MLflow was used to version and manage prompts and schemas, making agent behavior more reproducible and easier to audit. PostgreSQL was used to persist platform metadata such as agents, tools, workflows, and applications. Laminar was integrated to observe workflow execution, capture traces, monitor latency, follow token usage, and provide visibility into the behavior of agents during execution.

The implemented platform was validated through a concrete financial document analysis use case. The complete scenario demonstrated the main lifecycle of the platform: creating an intelligent application, configuring agents, selecting models and tools, associating prompts and schemas, deploying generated agents as Docker containers, composing a workflow, executing the workflow, and inspecting execution traces. This scenario confirmed the feasibility of the proposed architecture and showed how complex business processes can be transformed into modular and observable multi-agent workflows.

In parallel with VIP, the SLM-Evals project was developed as a complementary evaluation pipeline for comparing large and small language models. This second project focused on the evaluation of model outputs using a financial analysis case, a golden dataset, baseline Prompt Registration, MLflow experiment tracking, multi-layer evaluation metrics, and prompt enhancement experiments. The SLM-Evals pipeline provided a structured way to compare models such as Phi-4 and Llama 3.1 according to quality, structure, semantic alignment, numerical reliability, latency, and cost-related indicators.

The two parts of the internship are therefore complementary. VIP provides the operational platform for creating and orchestrating intelligent agents, while SLM-Evals provides an experimental framework for evaluating model behavior and supporting better provider selection. Together, they contribute to the same objective: making enterprise AI systems more modular, controlled, observable, and adaptable to business requirements.

Overall, this project allowed the development of a complete technical foundation for an agentic AI platform. It combined backend development, frontend integration, Docker-based deployment, multi-agent orchestration, schema validation, prompt management, observability, and model evaluation. The obtained result is not only an implementation prototype, but also a reusable architectural blueprint for future intelligent applications in financial and enterprise contexts.

\subsection{Current Limitations}

Although the implemented solution validates the main objectives of the project, several limitations remain and should be considered before extending the platform toward a production-grade environment.

The first limitation concerns the scope of validation. The VIP platform was mainly validated through one financial document intelligence workflow composed of extractor, mapper, generator, and sanitizer agents. This scenario demonstrates the core capabilities of the platform, but it does not yet prove the behavior of the architecture across a large number of business domains, document types, workflow structures, or user profiles. Additional use cases would be required to confirm the generalizability of the platform.

The second limitation is related to the SLM-Evals evaluation dataset. The evaluation pipeline was implemented and tested on a controlled financial analysis case. This was sufficient to validate the structure of the evaluation process, the MLflow tracking logic, and the comparison between model outputs. However, a larger and more diverse dataset would be necessary to draw stronger conclusions about model performance. Future evaluations should include more documents, different financial institutions, several reporting periods, multilingual inputs, and more heterogeneous document layouts.

Another limitation concerns the maturity of the deployment environment. The current implementation mainly relies on local Docker-based deployment and service-to-service communication through a shared Docker network. This validates the containerized execution model, but it does not yet include all the mechanisms required for production deployment, such as Kubernetes orchestration, autoscaling, high availability, centralized secrets management, load balancing, and environment separation between development, testing, and production.

The workflow orchestration layer also has some limitations. The current workflow execution is mostly sequential and based on predefined workflow definitions. This is suitable for the implemented scenario, but more advanced orchestration capabilities would be needed for complex enterprise workflows. These include conditional branching, parallel execution, retries, timeout policies, fallback agents, human validation steps, and workflow versioning.

The adapter mechanism is another important limitation. The adapter helps transform outputs between agents when schemas are not naturally aligned. However, when this transformation depends on an LLM, the behavior may not always be fully deterministic. For stable production workflows, deterministic mapping rules or validated transformation templates should be preferred whenever possible.

Governance features are also still incomplete. MLflow is used for prompt and schema versioning, but the platform does not yet include a complete governance workflow for approving, promoting, rejecting, or rolling back prompt and schema versions. Similarly, advanced access control, audit policies, role-based permissions, and validation gates should be strengthened before using the platform in sensitive enterprise environments.

Finally, the observability layer provides useful traces and execution details, but it can still be improved. Current observability focuses on workflow traces, latency, token usage, and execution inspection. Future versions should include alerting mechanisms, quality monitoring dashboards, failure classification, long-term execution analytics, and automatic detection of abnormal behavior.

\subsection{Perspectives for Future Work}

Several perspectives can extend and improve the work completed during this internship.

A first perspective is to expand the number of supported business use cases in VIP. The current financial document intelligence scenario can be extended to other financial workflows such as credit file analysis, risk report generation, compliance checking, client document summarization, and automated financial ratio interpretation. Testing the platform on several use cases would help validate the flexibility of the configuration-driven approach and demonstrate the reuse of generated agents across different business processes.

A second perspective is to strengthen the integration between VIP and SLM-Evals. The evaluation results produced by SLM-Evals could be used directly inside VIP to guide model selection for each agent. For example, an extractor agent may require strong factual precision, a generator agent may require high reasoning and writing quality, while a sanitizer agent may require strict format compliance. The platform could recommend the most suitable model according to task type, cost, latency, and quality requirements.

Another important improvement is the extension of the SLM-Evals dataset. Future experiments should include multiple golden datasets, more financial documents, multilingual examples, scanned and non-scanned files, and more complex tables. This would make the evaluation process more reliable and would allow better comparison between proprietary LLMs, open-source LLMs, and smaller language models. A larger benchmark would also help identify whether SLMs can replace larger models for specific agent roles.

The platform can also be improved by adding stronger governance mechanisms. Future versions should include role-based access control, user permissions, prompt approval workflows, schema validation gates, version promotion, automatic rollback, and detailed audit logs. These features are especially important for enterprise and financial environments where traceability, control, and compliance are required.

From a deployment perspective, the next step would be to move from local Docker-based execution to a more production-oriented infrastructure. Kubernetes could be introduced to manage container orchestration, service discovery, autoscaling, health checks, and rolling updates. A CI/CD pipeline could also be added to automate testing, image building, deployment, and environment promotion.

Workflow orchestration can also be enriched. Future versions of VIP could support conditional workflows, parallel execution, reusable workflow templates, retry policies, timeout management, and human-in-the-loop validation. This would allow the platform to support more realistic enterprise processes where some decisions require validation, exception handling, or dynamic routing.

Another perspective is to improve the Agent Factory. The generation mechanism could include automatic pre-deployment testing, schema compatibility checks, tool permission verification, dependency security scanning, and generated agent health validation. This would make the creation of new agents safer and more reliable.

The tool catalog can also be extended. Future work could add connectors for databases, APIs, document storage systems, financial calculators, OCR tools, table extraction modules, and business-specific services. A richer tool ecosystem would make generated agents more capable while preserving explicit control over what each agent is allowed to do.

Observability can be further developed into a complete monitoring and supervision layer. In addition to traces, the platform could include dashboards for workflow success rates, average latency, token consumption, model cost, error categories, schema validation failures, and agent-level performance. Alerting mechanisms could notify users when abnormal behavior or repeated failures are detected.

Finally, future work can focus on improving reliability and trust. This includes fallback models, deterministic post-processing, automatic output repair, stronger schema enforcement, confidence scoring, and evaluation before deployment. These mechanisms would help transform VIP from an experimental platform into a more robust enterprise-grade agentic AI hub.

In conclusion, the work completed during this internship provides a strong foundation for a modular, configurable, and observable AI agent platform. The next development steps should focus on scaling the number of use cases, strengthening governance, improving deployment maturity, and connecting model evaluation more directly to agent configuration and production decisions.

